% ============================================================
%  sections/05_results.tex
%  5. Kết quả & Thảo luận  (~2 trang)
% ============================================================

\section{Kết Quả và Thảo Luận}
\label{sec:results}

\subsection{Kết Quả Baseline Zero-Shot}

Bảng~\ref{tab:baseline} trình bày kết quả của \texttt{roberta-base-openai-detector}
trên bộ kiểm tra. Chỉ có tiếng Anh vượt mức ngẫu nhiên (64,1\%), trong
khi ba ngôn ngữ còn lại đều ở mức 48--51\% --- sát ngưỡng phân loại
ngẫu nhiên (50\%). Đặc biệt, tiếng Ả Rập đạt Macro-F1 chỉ
\textbf{0,336}, cho thấy mô hình đang phân loại sai hệ thống: dự đoán
``machine'' cho văn bản human và ngược lại.

\begin{table}[h]
\centering
\caption{Kết quả Phương pháp 1: Baseline Zero-Shot (roberta-base-openai-detector)}
\label{tab:baseline}
\begin{tabular}{lrrrl}
\toprule
\textbf{Ngôn ngữ} & \textbf{Accuracy} & \textbf{Macro-F1} & \textbf{N test} & \textbf{Nhận xét} \\
\midrule
EN & 64,1\% & 0,623 & 800 & Vượt ngẫu nhiên \\
VI & 50,7\% & 0,493 & 797 & Gần ngẫu nhiên \\
ZH & 47,6\% & 0,407 & 800 & Dưới ngẫu nhiên \\
AR & 48,3\% & 0,336 & 800 & Phân loại sai hệ thống \\
\bottomrule
\end{tabular}
\end{table}

Kết quả này xác nhận \textbf{thiên lệch Anh ngữ nghiêm trọng} trong
bộ phát hiện tiền huấn luyện: triển khai trực tiếp cho dữ liệu đa
ngôn ngữ không chỉ kém hiệu quả mà còn có thể gây hại hơn không dùng gì.
Đường cong ROC của cả bốn ngôn ngữ được thể hiện tại
Hình~\ref{fig:roc_baseline}.

\begin{figure}[h]
\centering
\includegraphics[width=0.75\textwidth]{roc_baseline.png}
\caption{Đường cong ROC --- Baseline (roberta-base-openai-detector).
Chỉ EN đạt AUC khả dụng (>0.6); AR, VI, ZH gần như ngẫu nhiên.}
\label{fig:roc_baseline}
\end{figure}

\subsection{Kết Quả Tinh Chỉnh XLM-RoBERTa}

Sau khi tinh chỉnh trên 12.789 mẫu đa ngôn ngữ, hiệu năng cải thiện
đáng kể trên tất cả bốn ngôn ngữ (Bảng~\ref{tab:finetuned}). Mức tăng
trung bình ước tính là \textbf{+28,7 điểm phần trăm} so với baseline.

\begin{table}[h]
\centering
\caption{Kết quả Phương pháp 2: Tinh Chỉnh XLM-RoBERTa (xlm-roberta-base)}
\label{tab:finetuned}
\begin{tabular}{lrrrr}
\toprule
\textbf{Ngôn ngữ} & \textbf{Accuracy} & \textbf{Macro-F1} & \textbf{N test}
    & \textbf{$\Delta$ Acc} \\
\midrule
EN & \textbf{95,4\%} & 0,954 & 800  & +31,3 pp \\
VI & 70,3\%          & 0,702 & 797  & +19,6 pp \\
ZH & 71,5\%          & 0,712 & 800  & +23,9 pp \\
AR & \textbf{88,3\%} & 0,882 & 800  & +40,0 pp \\
\midrule
\textbf{Tổng} & \textbf{81,4\%} & \textbf{0,814} & 3.197 & \textbf{+28,7 pp} \\
\bottomrule
\end{tabular}
\end{table}

Tiếng Ả Rập đạt mức tăng lớn nhất (+40,0 pp), từ trạng thái ``phân
loại sai hệ thống'' lên 88,3\% --- cho thấy XLM-RoBERTa đã học được
các đặc trưng hình thái của tiếng Ả Rập một cách hiệu quả. Tiếng Việt
và tiếng Trung đạt khoảng 70--72\%, thấp hơn đáng kể so với EN và AR,
phản ánh độ khó nội tại: Teencode trong tiếng Việt và văn bản Weibo cực
ngắn (15--30 ký tự) trong tiếng Trung. Hình~\ref{fig:roc_finetuned} hiển
thị đường cong ROC sau tinh chỉnh.

\begin{figure}[h]
\centering
\includegraphics[width=0.75\textwidth]{roc_finetuned.png}
\caption{Đường cong ROC --- XLM-RoBERTa sau tinh chỉnh.
EN và AR đạt AUC gần hoàn hảo (>0.95); VI và ZH cải thiện đáng kể so với baseline.}
\label{fig:roc_finetuned}
\end{figure}

\subsection{Kết Quả Phát Hiện Dựa Trên Perplexity}

Bảng~\ref{tab:ppl} và Hình~\ref{fig:roc} trình bày kết quả phương pháp
PPL. EN (AUC = 0,754) và AR (AUC = 0,778) đạt hiệu năng khá tốt, trong
khi VI (AUC = 0,538) và ZH (AUC = 0,534) gần như ngẫu nhiên.

\begin{table}[h]
\centering
\caption{Kết quả Phương pháp 3: Perplexity Zero-Shot (Qwen2.5-1.5B)}
\label{tab:ppl}
\begin{tabular}{lrrrrl}
\toprule
\textbf{Ngôn ngữ} & \textbf{AUC-ROC} & \textbf{PPL\textsubscript{Human}}
    & \textbf{PPL\textsubscript{Machine}} & \textbf{Tỉ lệ} & \textbf{Đánh giá} \\
\midrule
EN & \textbf{0,754} & 1.193,6 & 187,1  & 6,4$\times$ & Tốt \\
AR & \textbf{0,778} & 260,7   & 76,7   & 3,4$\times$ & Tốt \\
VI & 0,538          & 1.210,4 & 796,0  & 1,5$\times$ & Yếu \\
ZH & 0,534          & 221,5   & 207,9  & 1,07$\times$ & Ngẫu nhiên \\
\bottomrule
\end{tabular}
\end{table}

\begin{figure}[h]
\centering
\includegraphics[width=0.75\textwidth]{roc_statistical.png}
\caption{Đường cong ROC của phương pháp PPL (Qwen2.5-1.5B) trên 4 ngôn ngữ.
EN và AR có đường cong gần góc trên trái; ZH gần như trùng đường chéo ngẫu nhiên.}
\label{fig:roc}
\end{figure}

Hiệu năng kém của VI và ZH được lý giải qua tỉ lệ PPL: ZH chỉ đạt
\textbf{1,07$\times$} --- con người và máy hầu như không phân biệt được
qua PPL --- do bài đăng Weibo cực ngắn khiến cả hai lớp có entropy tương
đương nhau. VI đạt 1,5$\times$ nhưng chưa đủ lớn để phân tách tốt, do
Teencode làm tăng PPL đồng đều cho cả hai lớp.

\subsection{Kết Quả Học Chuyển Giao (Transfer Learning)}

Bảng~\ref{tab:transfer} trình bày kết quả của bốn thí nghiệm
\textit{leave-one-language-out}. Hiệu năng trung bình chỉ đạt
\textbf{54,9\% accuracy} và \textbf{0,543 F1}, gần với mức ngẫu nhiên.

\begin{table}[h]
\centering
\caption{Kết quả học chuyển giao (xlm-roberta-base, leave-one-language-out)}
\label{tab:transfer}
\begin{tabular}{lcccc}
\toprule
\textbf{Ngôn ngữ} & \textbf{Ngôn ngữ huấn luyện} & \textbf{Acc}
    & \textbf{F1} & \textbf{ROC-AUC} \\
\midrule
EN & VI, ZH, AR & 0,545 & 0,683 & 0,851 \\
VI & EN, ZH, AR & 0,531 & \textbf{0,134} & 0,536 \\
ZH & EN, VI, AR & 0,503 & 0,666 & 0,512 \\
AR & EN, VI, ZH & \textbf{0,619} & 0,691 & \textbf{0,810} \\
\midrule
\textbf{Trung bình} & & \textbf{0,549} & \textbf{0,543} & \textbf{0,677} \\
\bottomrule
\end{tabular}
\end{table}

Kết quả cho thấy ba hiện tượng đáng chú ý:

\begin{enumerate}
    \item \textbf{Tiếng Việt thất bại hoàn toàn} (F1 = 0,134, recall = 0,076):
    Mô hình dự đoán gần như toàn bộ là ``human'', cho thấy đặc trưng
    MGT tiếng Việt (Teencode, cấu trúc câu không chính thức) không thể
    học được từ các ngôn ngữ khác.

    \item \textbf{Tiếng Trung bị thoái hóa} (Acc = 50,3\%): Mô hình thiên
    về nhãn ``machine'' (395/400 human bị phân loại sai), dẫn đến recall
    cao giả tạo (0,99) nhưng precision ngang tung đồng (0,501).

    \item \textbf{EN và AR được hưởng lợi một phần}: AR đạt kết quả tốt
    nhất (Acc = 61,9\%, AUC = 0,81). EN đạt recall 0,98 nhưng precision
    chỉ 0,524 --- thiên lệch về nhãn ``machine''.
\end{enumerate}

\subsection{So Sánh Tổng Hợp Bốn Phương Pháp}

Bảng~\ref{tab:comparison} tổng hợp hiệu năng của bốn phương pháp.

\begin{table}[h]
\centering
\caption{So sánh tổng hợp bốn phương pháp}
\label{tab:comparison}
\begin{tabular}{lrrrr}
\toprule
\textbf{Ngôn ngữ} & \textbf{Baseline Acc} & \textbf{XLM-R Acc}
    & \textbf{PPL AUC} & \textbf{Transfer Acc} \\
\midrule
EN & 64,1\% & \textbf{95,4\%} & 0,754 & 54,5\% \\
VI & 50,7\% & \textbf{70,3\%} & 0,538 & 53,1\% \\
ZH & 47,6\% & \textbf{71,5\%} & 0,534 & 50,3\% \\
AR & 48,3\% & \textbf{88,3\%} & 0,778 & 61,9\% \\
\midrule
\textbf{Khả dụng} & Chỉ EN & \textbf{4 ngôn ngữ} & EN + AR & Không \\
\bottomrule
\end{tabular}
\end{table}

XLM-RoBERTa tinh chỉnh là phương pháp vượt trội nhất trên tất cả ngôn
ngữ và tiêu chí. Phương pháp PPL có ưu điểm zero-shot (không cần dữ
liệu nhãn) và hiệu quả với EN/AR, nhưng thất bại với VI/ZH. Baseline
tiếng Anh không thể dùng được cho bài toán đa ngôn ngữ. Học chuyển giao
cho thấy mỗi ngôn ngữ cần dữ liệu huấn luyện riêng --- đặc biệt với
các ngôn ngữ có đặc thù hình thái như tiếng Việt và tiếng Trung.
